% Benötigt die Farben klasseHell, klasseDunkel, klasseHellText aus classicthesis-config.tex
\begin{figure}[H]
\centering
\begin{tikzpicture}
\begin{axis}[
    name=vs, 
    width=\linewidth, height=4.8cm,
    ybar=2pt, bar width=14pt, enlarge x limits=0.2,
    symbolic x coords={HitRate@1,HitRate@5,Precision@5}, xtick=data,
    ymin=0, ymax=55, ylabel={in \%}, ymajorgrids=true,
    yticklabel style={/pgf/number format/.cd, use comma, 1000 sep={\,},},
    nodes near coords, nodes near coords style={font=\scriptsize, /pgf/number format/.cd, use comma, fixed, fixed zerofill, precision=1},
    title={Vorderseite (156 Anfragen)}, title style={font=\small},
    xtick pos=bottom, ytick pos=left,
    axis line style={gray!60}, tick style={gray!60},
    grid style={gray!25},
    ticklabel style={font=\small},
    label style={font=\small},
]
\addplot[fill=klasseHell, draw=none] coordinates {(HitRate@1,16.7) (HitRate@5,43.6) (Precision@5,11.0)};
\addplot[fill=klasseHellText, draw=none] coordinates {(HitRate@1,16.7) (HitRate@5,42.9) (Precision@5,10.9)};
\addplot[fill=klasseDunkel, draw=none] coordinates {(HitRate@1,17.3) (HitRate@5,46.2) (Precision@5,12.4)};
% Zufallswerte, sobald im Code berechnet (Werte in % eintragen und auskommentieren):
% \draw[gray!70!black, dashed] ({rel axis cs:0,0} |- {axis cs:HitRate@1,4.5}) -- ({rel axis cs:1,0} |- {axis cs:HitRate@1,4.5});
\end{axis}
\begin{axis}[
    name=rs, at={(vs.below south west)}, anchor=north west, yshift=-0.9cm,
    width=\linewidth, height=4.8cm,
    ybar=2pt, bar width=14pt, enlarge x limits=0.2,
    symbolic x coords={HitRate@1,HitRate@5,Precision@5}, xtick=data,
    ymin=0, ymax=55, ylabel={in \%}, ymajorgrids=true,
    yticklabel style={/pgf/number format/.cd, use comma, 1000 sep={\,},},
    nodes near coords, nodes near coords style={font=\scriptsize, /pgf/number format/.cd, use comma, fixed, fixed zerofill, precision=1},
    title={Rückseite (62 Anfragen)}, title style={font=\small},
    xtick pos=bottom, ytick pos=left,
    axis line style={gray!60}, tick style={gray!60},
    grid style={gray!25},
    ticklabel style={font=\small},
    label style={font=\small},
    legend style={at={(0.5,-0.25)}, anchor=north, legend columns=3, font=\footnotesize, draw=none, /tikz/every even column/.append style={column sep=8pt}},
    legend image code/.code={\draw[#1, draw=none] (0cm,-0.1cm) rectangle (0.3cm,0.1cm);},
]
\addplot[fill=klasseHell, draw=none] coordinates {(HitRate@1,14.5) (HitRate@5,29.0) (Precision@5,9.4)};
\addplot[fill=klasseHellText, draw=none] coordinates {(HitRate@1,14.5) (HitRate@5,29.0) (Precision@5,9.4)};
\addplot[fill=klasseDunkel, draw=none] coordinates {(HitRate@1,16.1) (HitRate@5,27.4) (Precision@5,8.7)};
\legend{Kosinus, Euklidisch, Manhattan}
% Zufallswerte, sobald im Code berechnet (Werte in % eintragen und auskommentieren):
% \draw[gray!70!black, dashed] ({rel axis cs:0,0} |- {axis cs:HitRate@1,4.5}) -- ({rel axis cs:1,0} |- {axis cs:HitRate@1,4.5});
\end{axis}
\end{tikzpicture}
\caption[Kennzahlen der drei Ähnlichkeitsmaße]{HitRate@1, HitRate@5 und Precision@5 der drei Ähnlichkeitsmaße auf Vorder- und Rückseite.}
\label{fig:distanzvergleich}
\end{figure}
\todo{Balken Werte Zufall eintragen}
